\documentclass{article}
\usepackage{microtype}
\usepackage{booktabs}
\usepackage{hyperref}
\usepackage{icml2026}
\usepackage{amsmath}

\icmltitlerunning{Can a Discriminative Judge Recover Credit from Failed Agent Traces?}

\begin{document}
\twocolumn[
\icmltitle{Can a Discriminative Judge Recover Credit from Failed Agent Traces?}
\begin{icmlauthorlist}
\icmlauthor{Anonymous Author}{anon}
\end{icmlauthorlist}
\icmlaffiliation{anon}{Anonymous Institution}
\icmlcorrespondingauthor{Anonymous Author}{anonymous@example.com}
\vskip 0.3in
]
\printAffiliationsAndNotice{}

\begin{abstract}
Long-horizon agents often receive only an outcome reward after many tool calls, making it difficult to distinguish productive actions from the mistake that ultimately causes failure. We study whether a frozen, typed discriminative judge can supply useful transition-level measurements without training a task-specific process reward model. Our prefix-only protocol exposes public observations, actions, and static success criteria while withholding verifier outputs and oracle fields. Matched frozen-trace comparisons give heterogeneous evidence: Jev outperforms a local 8B language-model judge on ScienceWorld failure-step localization across three policy sizes, while naive score aggregation underperforms random best-of-five selection on two panels and ToolSandbox's hidden milestones favor the language-model judge. We then test the signal in matched ALFWorld policy training. On an identical first batch where all 16 trajectories fail, sparse GRPO gives zero advantage while prefix-only Jev gives nonzero turn-specific advantage to all 480 transitions. After ten updates, held-out success is 19/64 for prefix-only Jev versus 1/64 for sparse GRPO. A single matched post-episode arm, which conditions Jev on the completed public trace and verified outcome, reaches 22/64; its final mean verifier score is nevertheless lower than prefix-only Jev's (1.250 versus 1.423). These one-seed results motivate replication but do not establish general sample efficiency or uniform superiority.
\end{abstract}

\section{Introduction}

An agent may inspect a relevant object, call several useful tools, and then take one irreversible wrong action. If training receives only a terminal success bit, the entire failed rollout is labeled alike. This ambiguity becomes acute when all samples in a Group Relative Policy Optimization (GRPO) group fail: a group-normalized binary outcome reward has no between-rollout variance, even though the actions may have very different effects on the environment \citep{shao2024deepseekmath,tao2026trace}. Outcome verification remains valuable because it checks the actual goal; the problem is that a final bit alone does not identify where the trajectory became better or worse.

Process supervision offers an attractive remedy, but obtaining reliable step labels is itself difficult. In mathematical reasoning, human-annotated process rewards have proved useful \citep{lightman2023verify}; interactive agents add state changes, tool side effects, and partially observed goals, for which a locally plausible action need not be globally useful \citep{choudhury2025agentprm,xi2025agentprm}. Native environment scores provide a reference when available, but many tool-use tasks have no inexpensive step verifier. A generative language-model judge can interpret such traces, yet even a JSON-constrained non-thinking judge incurs autoregressive output and parsing cost. We therefore ask whether a frozen discriminative judge can measure task-relevant transition quality directly from the public state, action, and observed result. Jev returns typed Score, Choice, and Noul decisions with numeric outputs rather than generated explanations \citep{typesafe2026docs}; this is a design affordance, not evidence that its judgments are accurate or cheap on agent traces.

We define the primary target problem as \emph{credit recovery among outcome-equivalent failures}. Given a task $x$, a rollout $\tau=(o_0,a_0,o_1,\ldots,o_T)$, and terminal verifier $R(\tau)$, can a prefix judge $j_t=f(x,o_{\leq t},a_{\leq t},o_{t+1})$ identify useful progress and the first damaging transition when $R(\tau)=0$ for all rollouts in a group? This judge sees only information available by transition $t$; it never sees later actions, terminal reward, hidden oracle state, or gold paths. We separately study a hindsight annotator $j_t^{H}=f(x,\tau,R(\tau),t)$ that labels every transition after the episode ends. Future and outcome information are legitimate inputs to post-episode RL labels because the deployed policy does not consume the annotator, but this second estimator is noncausal and cannot support a prefix-PRM claim. Both formulations distinguish temporal localization from merely ranking complete rollouts by a scalar, and both make density an insufficient success criterion: a constant or confidently wrong step score is dense but provides no valid credit.

Our frozen evaluation panel contains 600 Qwen3 trajectories and 5,743 transitions across ToolSandbox and ScienceWorld, with five stochastic rollouts per task, three policy sizes, and 20 tasks per benchmark/model pair. These environments expose different evaluation anchors: ToolSandbox records stateful tool outcomes and milestones \citep{lu2024toolsandbox}, whereas ScienceWorld provides native score changes in a text-interactive environment \citep{wang2022scienceworld}. In the Qwen3-8B panel, all five rollouts fail in 16 of 20 ToolSandbox groups and 17 of 20 ScienceWorld groups. These counts establish a substantial zero-variance slice for \emph{binary} outcome GRPO; they do not establish that Jev recovers the missing signal. Continuous terminal verifier scores may still distinguish failed trajectories and must be retained as a stronger baseline.

The study has three linked tests. First, on frozen traces, we compare prefix-only Jev with binary and continuous outcome rewards, native process anchors, and a matched language-model judge. We evaluate failure-only progress discrimination, first-fault localization, best-of-five verifier utility, calibration, and quality-adjusted latency and cost. Second, we train matched policies with the same GRPO implementation and benchmark split, changing only the advantage source, then compare held-out success at matched environment steps and compute. Third, without rerunning the first two arms, we add one otherwise matched outcome-conditioned hindsight arm. Offline evidence can support only the claim that a judge \emph{measures} process credit better; a downstream improvement requires policy training. All training comparisons currently use one seed, so the broader sample-efficiency claim remains contingent on replication.

\section{Related Work}

\paragraph{Outcome-based agentic reinforcement learning and delayed credit.}
GRPO replaces a learned value baseline with within-group reward comparisons \citep{shao2024deepseekmath}. Its simple outcome formulation is useful when a verifier is trustworthy, but equal rewards produce no group-relative contrast. TRACE studies the separate long-horizon problem of assigning turn-level credit from a frozen reference model and reports training gains in search agents \citep{tao2026trace}. Our judge-based question differs: we ask whether a \emph{general-purpose frozen discriminator}, without a task-specific credit estimator or future outcome information, can recover verifiable transition distinctions among failed tool trajectories. We preserve continuous terminal rewards as an outcome-only control so that any benefit cannot be attributed merely to replacing a success bit with a more graded final score.

\paragraph{Process reward models and agent evaluation.}
Process supervision was shown to improve verification on mathematical solutions with human step annotations \citep{lightman2023verify}. For agents, AgentPRM and related methods learn progress-sensitive signals from interaction or demonstrations \citep{choudhury2025agentprm,xi2025agentprm}. ToolPRMBench emphasizes local error discrimination in realistic tool-use histories and shows why intermediate judgment needs its own evaluation rather than being inferred from final-answer accuracy \citep{li2026toolprmbench}. Our setting differs in using Jev as a zero-shot annotator: no PRM parameters are trained before scoring the frozen traces. This removes a PRM-training stage but does not remove the need to validate labels, guard against reward hacking, or test generalization during policy optimization.

\paragraph{Discriminative judging and measurable process credit.}
Jev accepts a state and independently evaluates typed questions, returning a distribution over descriptive Score levels or a Noul probability \citep{typesafe2026docs}. We use this interface for a narrowly specified transition-effect judgment and retain the full response distribution. The vendor documents limitations on numerical precision, irrelevant context, and adversarial text \citep{typesafe2026jaggedness}; consequently, arithmetic is done outside the model and inputs are filtered to task-relevant public fields. Unlike a generic ``rate this trajectory'' prompt, this yields a falsifiable local signal whose agreement with hidden verifier anchors, temporal localization, latency, and downstream GRPO effects can be measured separately.

\section{Methods}

For each transition, the judge input is a whitelist of the user task, an explicit static specification of every condition required for success, at most three earlier public action--observation pairs, the current public observation, the action, and the observed tool or environment response. It excludes the verifier's judgment on the current trajectory, process rewards, oracle diagnostics, gold action sequences, and every future step. Revealing the success predicate is necessary to define process quality; revealing its per-trajectory output would leak the label. All offline judges score the same previously sampled trajectories, already labeled by the benchmark verifier; no second policy rollout is generated. Judge records are matched to frozen source records by trajectory ID and step index, with source checksums retained. The frozen comparisons below retain the original three-level rubric for reproducibility. The revised online rubric instead has harmful and beneficial endpoints and uses Jev's native continuous expected Score $q_{i,t}\in[0,1]$ together with its native confidence $c_{i,t}\in[0,1]$; neither value is discretized.

The offline primary slice comprises groups whose five binary outcome rewards are all zero. We compare step labels against held-out native process anchors, report failure-only AUROC and first-fault error, and measure best-of-five terminal verifier utility without exposing those anchors to any judge. Quality must be reported jointly with throughput and token cost. For online training, the sparse RLVR baseline uses $R_i=10\,\mathbf{1}[\mathrm{success}_i]$, standardizes the four trajectory returns for a task, and broadcasts the resulting $A_i^{\mathrm{out}}$ to every action in trajectory $i$. The revised Jev arm assigns the direct turn-level advantage
\begin{equation}
\begin{aligned}
A_{i,t}^{\mathrm{Jev}}&=c_{i,t}(2q_{i,t}-1),\\
A_{i,t}&=A_{i,t}^{\mathrm{Jev}}+\beta A_i^{\mathrm{out}},\qquad \beta=0.1.
\end{aligned}
\end{equation}
The Jev term is not group-standardized: its absolute magnitude encodes confidence, and standardization would amplify nearly neutral judgments. Only action tokens emitted at turn $t$ receive $A_{i,t}$. With token-level importance ratio $\rho_{i,t,k}$, the actor minimizes the usual clipped PPO surrogate $-\min(\rho A_{i,t},\operatorname{clip}(\rho,1-\epsilon,1+\epsilon)A_{i,t})$. In the prefix arm, Jev is queried after each newly generated transition and receives no later state or terminal result. In the hindsight arm, one request is issued after termination with the complete public trajectory, verified reward and success bit, and one typed Score question per transition; it uses the same $A_{i,t}$ equation. Oracle state, gold paths, stored process rewards, and hidden verifier internals remain excluded. Neither arm reads frozen annotations, and held-out evaluation uses only the environment verifier. All controlled runs start from Qwen2.5-1.5B-Instruct with seed 0, four prompts and four rollouts per update, a 30-action horizon, ten updates, and \texttt{valid\_seen} evaluations at updates 0, 5, and 10. The earlier aggregate and reward-to-go variants are retained only as ablations. These one-seed pilots are not a general sample-efficiency test.

\section{Experiments}

Table~\ref{tab:frozen-summary} summarizes the matched frozen-trace experiments before the detailed results. AUROC is computed within a task group or game and then averaged so that long trajectories do not dominate. The number of eligible units is reported because several panels contain many all-failure groups but few groups with both positive and negative native process events.

\begin{table*}[t]
\caption{Summary of frozen-trace process-credit measurements. Each Jev/LLM pair scores the same public states and semantic criterion. ``Groups'' denotes eligible task groups or games for the reported equal-weighted AUROC.}
\label{tab:frozen-summary}
\centering
\small
\begin{tabular}{llrrrr}
\toprule
Panel & Process target & Steps & Groups & Jev & Local 8B LLM \\
\midrule
ToolSandbox, Qwen3-8B & positive official milestone delta & 215 & 3 & 0.353 & 0.720 \\
ScienceWorld, Qwen3-8B & positive native score delta & 2,069 & 10 & 0.876 & 0.502 \\
ScienceWorld, Qwen3-4B & positive native score delta & 1,578 & 7 & 0.865 & 0.560 \\
ScienceWorld, Qwen3-1.7B & positive native score delta & 1,596 & 3 & 0.989 & 0.388 \\
ALFWorld dev, Qwen3-8B & shortest-plan progress & 306 & 5 & 0.892 & 0.797 \\
ALFWorld dev, Qwen3-8B & hard-step quality & 306 & 6 & 0.773 & 0.660 \\
ALFWorld failures, Qwen3-1.7B & shortest-plan progress & 300 & 4 & 0.563 & 0.628 \\
ALFWorld failures, Qwen3-1.7B & hard-step quality & 300 & 5 & 0.761 & 0.551 \\
\bottomrule
\end{tabular}
\end{table*}

The first complete judge panel covers all 215 transitions of 100 frozen Qwen3-8B ToolSandbox rollouts. In its 16 all-failed groups, every one of the 80 trajectories receives a nonzero Jev-based trajectory RLOO advantage, whereas binary outcome GRPO receives zero. Yet density alone fails the quality check: on nonterminal transitions, the mean within-trajectory AUROC for detecting positive official milestone deltas is 0.304 for Jev across only 14 eligible traces (three task groups); the corresponding within-group AUROC is 0.353. A matched local, non-thinking Qwen3-8B language-model judge, given the same state and rubric, obtains 0.679 and 0.720, respectively. Observed mean latency is 0.908 seconds per remote Jev call versus 1.195 seconds per local LLM call; this is not a hardware-matched speed comparison. These small, correlated strata do not support a general performance claim. The mismatch is informative: some hidden ToolSandbox milestones give credit to an action whose public tool response reports an error or no visible state change. A judge restricted to public prefixes cannot observe that hidden update. We keep this negative result rather than revise the rubric on the same test groups.

The complete paired pass over 2,069 previously sampled Qwen3-8B ScienceWorld transitions shows a different pattern. Among 17 all-five-failed groups, mean AUROC for positive native score changes is 0.900 for Jev within 45 eligible trajectories versus 0.490 for the matched non-thinking Qwen3-8B language-model judge. Within ten eligible task groups, the means are 0.876 and 0.502, respectively; Jev is higher in nine groups. The paired group-mean difference is 0.374, with an exploratory task-bootstrap 95\% interval of [0.236, 0.507] from 20,000 resamples. Both judges received identical public transition states and semantic criteria. An exploratory slice excluding repeated state--action pairs retains a mean group AUROC of 0.876 versus 0.416 over 627 steps, so the gap is not solely loop detection. Binary and continuous terminal rewards cannot localize different steps inside a fixed trajectory (AUROC 0.500 by construction), although continuous final reward remains a cross-trajectory control. Jev yields nonzero trajectory-level RLOO advantage on 83 of the 85 failed-group trajectories and varying reward-to-go on 64. This establishes a process-measurement advantage over this particular 8B judge on the eligible ScienceWorld tasks, not over all judges or in downstream learning.

A second complete pass on 1,578 frozen Qwen3-4B ScienceWorld transitions replicates the local-credit direction across policy sizes, with the same judge and rubric. In 16 all-five-failed groups, seven task groups have discriminable native score-change events; their mean within-group AUROC is 0.865 for Jev versus 0.560 for the LLM, with Jev higher in six groups. The paired mean difference is 0.305 and an exploratory task-bootstrap 95\% interval is [0.126, 0.493]. Across 25 eligible traces, mean within-trace AUROC is 0.809 versus 0.482. Jev has higher nonzero stepwise RLOO density among failed trajectories (0.988 versus 0.579), but density is not substituted for the AUROC result. Mean observed remote Jev latency is 1.015 s versus 0.848 s for the local LLM; hardware and network are unmatched. The policy panels share benchmark tasks, so they are not independent task replications.

A third complete pass covers 1,596 frozen Qwen3-1.7B ScienceWorld transitions. Nineteen of 20 five-rollout groups all fail, so binary outcome GRPO assigns zero group advantage to all 95 trajectories in that slice. Only three groups contain both positive and non-positive transitions. Mean within-group AUROC for detecting positive native score changes on those three groups is 0.989 for Jev versus 0.388 for the LLM; the paired difference is 0.601 with an exploratory task-bootstrap 95\% interval of [0.427, 0.867]. Jev wins all three groups. Across 15 eligible traces the corresponding means are also 0.989 and 0.388. Failure-only nonzero stepwise RLOO density is 0.991 versus 0.346. An exploratory control removes repeated state--action pairs, leaving 298 failed-group steps; Jev remains higher at 0.882 versus 0.262, so loop detection alone does not explain the gap. The near-ceiling primary result is strong local evidence but has narrow task support. Moreover, all sampled failed rollouts within each group have equal continuous verifier reward, so this panel cannot test best-of-five ranking utility.

The predeclared best-of-five test sharply limits the reward claim. On the 17 all-failed 8B-policy groups, selecting a rollout by the \emph{sum} of Jev transition scores yields mean official continuous final reward 0.079, below random selection (0.120), the matched LLM judge (0.129), and the sampled oracle maximum (0.151). Averaging Jev scores also yields 0.079. The 4B-policy panel repeats this negative trajectory-level result: Jev selects 0.010 under either sum or mean aggregation, below random 0.0196 and sampled oracle 0.0394. Thus local progress discrimination does not justify naive additive trajectory reward; the signed Jev signal may penalize long unsuccessful exploration more than the official partial-credit verifier does. Any training reward transformation must be chosen on development tasks and assessed on held-out tasks, not tuned to reverse this observed test result.

An exploratory ALFWorld development slice confirms the aggregation issue without establishing a training gain. Complete Jev annotations of 306 transitions from 12 frozen Qwen3-8B traces across six \texttt{valid\_unseen} games yield mean within-game AUROC 0.892 for shortest-plan progress (five eligible games) and 0.773 for hard-step quality (six games); a matched local non-thinking Qwen3-8B judge scores the exact same public states and obtains 0.797 and 0.660. A paired game-bootstrap gives exploratory 95\% intervals of [-0.007, 0.220] and [0.031, 0.209] for the two AUROC differences, respectively. First-hard-error exact localization is 6/12 for Jev versus 1/12 for this LLM. Mean observed call latency is 1.136 s for remote Jev and 0.871 s for the local LLM, which does not establish a hardware-matched speed advantage. In best-of-two selection, the mean centered Jev and LLM scores each pick a successful trace in three of six games, the sampled oracle ceiling, whereas the sum of Jev scores picks one of six and random selection has expectation two of six. The slice is too small for a general claim and cannot be reused as independent downstream test data after informing the reward aggregation choice.

An earlier, independent ALFWorld hard-failure set provides a contrasting diagnostic. On the same 300 frozen transitions from ten failed Qwen3-1.7B trajectories across five games, pooled failure-only shortest-plan-progress AUROC is 0.500 for terminal RLVR, 0.650 for Jev, and 0.629 for the matched LLM judge. Yet equal-weighted within-game AUROC is 0.500, 0.563, and 0.628 across four eligible games: the apparent Jev advantage over the LLM does not survive task balancing. For hard-step quality, equal-weighted within-game AUROC is 0.500, 0.761, and 0.551 across all five games, with Jev above the LLM in each game. Jev identifies the first hard error exactly in three of ten traces, versus zero for both comparators. The set is too small and correlated for a general superiority claim; these are offline measurement results, not evidence of improved GRPO learning.

The earlier trajectory-aggregate ALFWorld ablation starts both Qwen2.5-1.5B-Instruct arms at 0/64 held-out \texttt{valid\_seen} success. All 16 trajectories fail in each of the first five training groups. After removing reward fields and random identifiers, the arms' first 480 transition records have the same SHA-256 digest. The binary-outcome arm consequently has zero reward, zero advantage, and policy-gradient loss 0.000 at every update through five. The aggregate Jev arm yields nonzero trajectory advantages; at update 5, held-out success is 2/64 versus 1/64, and mean continuous verifier score is 0.101 versus 0.037. This one-success difference is inconclusive and the ablation does not assign distinct step advantages.

The revised confidence-preserving run begins from 0/64 held-out success in both arms. The two arms' first batches contain the same 480 public transitions after excluding run identifiers, and all 16 trajectories fail. Sparse GRPO therefore assigns zero advantage to every token and reports policy-gradient loss 0.000. Jev assigns nonzero advantage to 480/480 transitions, all 16 trajectories have nonconstant within-trajectory advantage, and the policy-gradient loss is 0.571. Advantages range from $-0.893$ to $0.951$ with mean absolute value 0.668. This establishes that distinct $A_{i,t}$ values reach the PPO objective rather than merely appearing in an annotation log. Over all ten updates, Jev's nonzero transition coverage is 0.986--1.000 and its nonconstant-trajectory fraction remains 1.000. Sparse GRPO receives an active outcome signal in only one of ten updates.

\begin{table*}[t]
\caption{Matched held-out ALFWorld \texttt{valid\_seen} evaluation. All arms start from the same checkpoint, use seed 0, and evaluate the same 64 tasks. Only the hindsight arm is newly run; earlier arms remain frozen.}
\label{tab:online-grpo}
\centering
\small
\begin{tabular}{lrrr}
\toprule
Evaluation & Sparse & Prefix Jev & Hindsight Jev \\
\midrule
Before training, success & 0/64 & 0/64 & 0/64 \\
Update 5, success & 0/64 & 8/64 & 16/64 \\
Update 5, mean score & 0.000 & 0.427 & 1.214 \\
Update 10, success & 1/64 & 19/64 & 22/64 \\
Update 10, mean score & 0.037 & 1.423 & 1.250 \\
\bottomrule
\end{tabular}
\end{table*}

The full prefix-only log contains 4,327 Jev calls, each with a continuous effect score and confidence and none with an API error. The request schema contains only the task, static success criteria, earlier public prefix, and current transition; recursive inspection finds no verifier output, stored reward, oracle field, gold path, or future step. At update 10, held-out success is 29.7\% for prefix-only Jev versus 1.6\% for sparse GRPO (Table~\ref{tab:online-grpo}). This satisfies the predeclared pilot gates---supervision under all-failure outcomes, dense turn-specific advantages that enter PPO, and a better trained-policy benchmark result. The magnitude justifies a formal multi-seed study, but one seed and 64 held-out tasks cannot establish a general Jev advantage, variance across training seeds, or sample-efficiency scaling.

The outcome-conditioned arm provides a controlled follow-up to the prefix-only estimator. Their first batches contain the same 480 public transitions under a canonical digest and all 16 trajectories fail. Hindsight Jev gives 479/480 transitions nonzero advantage, every trajectory has nonconstant within-trajectory credit, mean absolute advantage is 0.803, policy-gradient loss is 0.726, and gradient norm is 10.326. The corresponding prefix-only values on the same behavior are 480/480, 1.000, 0.668, 0.571, and 7.160. Thus outcome conditioning changes the magnitude and assignment of credit rather than merely making an already dense signal denser.

Across ten updates, hindsight Jev returns 4,350 transition labels in 160 completed-trajectory requests; 4,332 advantages are nonzero (99.6\%), every update has a nonconstant-trajectory fraction of 1.000, and no API error occurs. Mean and median request latency are 3.303 s and 1.755 s. The prefix and hindsight training-batch success counts are respectively $[0,0,3,3,4,1,0,2,6,11]$ and $[0,0,4,1,5,2,3,1,10,9]$ out of 16. These are descriptive on-policy batches, not held-out evaluations. The hindsight request count is about 27 times lower than the prefix arm's one-request-per-transition design, but each request contains a larger trace and many questions; call count alone is not a controlled throughput or cost comparison.

Strict held-out success favors hindsight at both checkpoints: 16/64 versus 8/64 at update 5 and 22/64 versus 19/64 at update 10. The final mean verifier score moves in the opposite direction, 1.250 versus 1.423. We therefore interpret the follow-up as evidence of faster early acquisition and a modest endpoint success gain at this seed, not uniform dominance. It supports testing outcome-conditioned labels across seeds while retaining the prefix-only arm as the causal measurement setting.

The superseded reward-to-go implementation was audited separately. It also produced dense turn-level advantages in an all-failure batch, but same-turn standardization erased the intended confidence scale. Its update-5 held-out success was 0/64 versus 1/64 for sparse GRPO, so it is retained only as an ablation and is excluded from Table~\ref{tab:online-grpo}.

\section{Discussion and limitations}

\paragraph{What the experiments establish.}
The most direct result is mechanical but important: when all 16 sampled trajectories fail, sparse binary GRPO has no within-group variance, while both Jev estimators assign nonconstant turn-level advantages that reach the clipped PPO loss. This resolves the signal-availability part of the delayed-credit problem. The frozen ScienceWorld panels additionally show that the prefix-only signal can align with native positive score changes substantially better than one matched local 8B judge across three policy sizes. Finally, the prefix-only training arm outperforms sparse GRPO on the same 64 held-out ALFWorld tasks after both five and ten updates. These claims concern the tested model, rubric, tasks, and seed.

\paragraph{What the experiments do not establish.}
Density is not accuracy. ToolSandbox reverses the ScienceWorld judge ordering because its hidden milestones are not always inferable from public observations, and naive sums or means of Jev scores underperform random best-of-five selection on two ScienceWorld panels. The hindsight arm improves strict ALFWorld success relative to prefix-only Jev at both checkpoints, but its final continuous verifier score is lower. The available evidence therefore rejects a universal ``Jev is a better reward'' claim. A more precise claim is that Jev can provide useful local credit when the success predicate and observable transition expose the relevant state change; the mapping from local judgments to an optimization reward remains task- and conditioning-dependent.

\paragraph{Causal prefix credit versus hindsight labels.}
The prefix arm tests whether a transition can be judged using only information available at that point. The hindsight arm instead asks a different and often easier question: knowing the complete public trajectory and verified outcome, which earlier transitions deserve positive or negative policy gradient? Using future information in a post-episode label is not a policy inference mismatch---RL returns and retrospective credit estimators do the same---but it prevents claims about real-time scoring or causal prefix quality. We therefore report the two arms separately and retain the prefix-only frozen benchmarks as the measurement test.

\paragraph{Statistical and systems limits.}
All policy-training results use one seed, ten updates, one small policy checkpoint, and 64 held-out \texttt{valid\_seen} tasks. The three ScienceWorld policy-size panels reuse benchmark tasks and are not independent task replications. The Jev service is remote and closed, while the comparison LLM is local; observed latency is neither hardware-matched nor a normalized dollar comparison. The hindsight request reduction counts HTTP calls rather than equal-sized units of work. Formal claims about sample efficiency, cost, or general benchmark superiority require multiple training seeds, broader task families, paired per-task uncertainty, and controlled serving measurements.

\section{Conclusion}

Across frozen traces and online training, Jev is most promising as a source of dense signed transition credit when sparse outcome groups collapse to a common reward. Prefix-only labels recover meaningful native process events on ScienceWorld and produce a large one-seed ALFWorld gain over sparse GRPO; outcome-conditioned hindsight labels retain the dense optimizer signal and reach a higher strict-success curve than prefix-only labels in the matched follow-up. The evidence is deliberately mixed: hidden-state ToolSandbox labels and trajectory-level selection expose failure modes, and the hindsight arm does not improve every held-out metric. The next decisive experiment is therefore multi-seed replication of the three training arms, not further post-hoc rubric tuning on the current tasks.

\bibliographystyle{icml2026}
\bibliography{references}
\end{document}
